% spec, 2026-10-24. Section "Related work". Source notes: agents/spec/prior_work.md (primary sources in agents/spec/ext/).
\section{Related work}\label{sec:related}

\paragraph{AlphaDev.} See Section~\ref{sec:background}. AlphaDev's only optimality statement is the brute-force claim for sort3
(about $10^{32}$ programs, with pruning, no certificate, search space not stated)~\cite{alphadev}. The kernels went into
LLVM libc++ (review D118029, landed 2022-04-08 as commit \texttt{194d1965d2c8})~\cite{d118029}, where they were accepted on
benchmarks; the patch makes no optimality claim.

\paragraph{Ullrich and Hack.} The closest work is ``Synthesis of Sorting Kernels''~\cite{ullrich} (CGO 2025). They search
sorting kernels with A$^\ast$/Dijkstra over abstract states, with optimality-preserving heuristics and optional
non-optimality-preserving cuts. They enumerate all optimal kernels for $n = 3$ (length 11 without memory instructions,
5602 programs after symmetry reduction) and report that no kernel of length 19 exists for $n = 4$ (exhaustive, two weeks), a
``new tight lower bound''; for $n = 5$ they find kernels of length about 33 with a cut, and give no lower bound.
Table~\ref{tab:uh} compares their model with ours, point by point. Each point where their model is \emph{more restrictive}
than AlphaDev's (fewer programs) means that their lower bound excludes fewer programs and so does not imply a lower bound
for AlphaDev's setting.

\begin{table}[h]\centering\small
\begin{tabular}{p{2.6cm}p{5.2cm}p{5.2cm}}\toprule
 & Ullrich and Hack~\cite{ullrich} & this work (MM via RS) \\\midrule
instructions & \texttt{mov}, \texttt{cmp}, \texttt{cmovl}, \texttt{cmovg} ($I_0$) & $I_1$ (adds \texttt{cmovle}, \texttt{cmovge}), as AlphaDev; $I_0$ as a side result \\
memory, counting & none; registers only (add $2n$ to compare) & loads and stores counted; exact map $L = 2n + k$ (Lemma~\ref{lem:snapshot}) \\
outputs & fixed in $r_1..r_n$ at the end (RF) & any register at any step (RS) = AlphaDev's stores \\
scratch registers & a fixed number $m$; the paper does not state $m$; the artifact code uses one (\texttt{SWAPS = 1}) & unbounded ($R = n + k$ suffices, Lemma~L3) \\
inputs tested & the $n!$ permutations (no ties) & all weak orders (ties); permutation subsets only inside sound relaxations \\
initial state & registers hold the inputs; reads of uninitialised scratch registers are forbidden & strict: garbage registers and flags (reads of uninitialised registers allowed); lower bound via a sound relaxation \\
pruning rules & argued in prose & each rule with a paper proof (and Lean where marked) \\
evidence & exhaustive-search runs, no proof log & DRAT/LRAT certificates checked by \texttt{cake\_lpr}; verified checkers in Lean and Rocq; two independent search tools \\
results & $n=3$: 11; $n=4$: 20; $n=5$: none & $n=3$: 11 (certified); $n=4$: 20 (exhaustive search, certificate checked to $k \le 17$); $n=5$: 32 best known, $\le 31$ open \\\bottomrule
\end{tabular}
\caption{Ullrich and Hack's model versus ours. In their model a lower bound of 11 / 20 does not imply one for AlphaDev's
setting: $I_0 \subsetneq I_1$, RF is stricter than RS, one scratch register is stricter than unbounded registers, and
forbidding reads of uninitialised registers excludes programs that are correct on x86 (e.g.\ a 0-\texttt{mov} sort3 core with one
such read, found by our search).}
\label{tab:uh}
\end{table}

\noindent Their $n = 3$ and $n = 4$ values agree with ours. Their lower bound transfers to tie-correct kernels in $I_0$ with one
scratch register, fixed outputs and no reads of uninitialised registers (distinct-input correctness is weaker than tie correctness), but not to $I_1$, to more
scratch registers, or to snapshot outputs. They also report that SMT, CEGIS, SyGuS (cvc5), ILP, PDDL planners and STOKE do not
solve $n = 4$, and that STOKE fails already for $n = 3$.

\paragraph{Shorter kernels outside the model.} Neri~\cite{neri} gives a 15-instruction branchless sort3 and a 14-instruction
sort3 with a loop; both use \texttt{sbb}/\texttt{adc}, table lookups and indexed stores, so they are outside the
\texttt{mov}/\texttt{cmp}/\texttt{cmov} model. Tunney~\cite{tunney} explains the AlphaDev saving as the deletion of one
\texttt{mov} from three compare-and-swap blocks (her 17-instruction x86 kernel passes our judge; it needs \texttt{cmovge}) and
shows a 14-instruction ARM64 kernel with the three-operand \texttt{csel}, also outside the model. Mimicry~\cite{mimicry} computes
a shuffle index instead of swapping: a 9-instruction SIMD sort3, and scalar kernels with indexed stores (sort3 17, sort4 27
for signed values); these use instructions outside our model and make speed claims only. None of them claims or refutes
optimality in AlphaDev's model.

\paragraph{Superoptimization.} Massalin's superoptimizer~\cite{massalin}, STOKE~\cite{stoke}, Lens~\cite{lens} and
Souper~\cite{souper} find short or fast instruction sequences by enumeration, stochastic search or solvers; none of them reports
an optimal sorting kernel, and STOKE fails on sort3 in the experiments of~\cite{ullrich}.

\paragraph{Optimal sorting networks.} The comparator model is solved for small $n$ with SAT and computer-checked certificates:
9 inputs need 25 comparators~\cite{codish}, with a Coq-verified checker of the certificate~\cite{cruzfilipe}, and 11 and 12
inputs need 35 and 39 comparators, with an Isabelle-verified checker~\cite{harder}. These results count comparators, not
instructions, and the 0-1 principle holds there. For cmov code a comparator is not atomic: one \texttt{cmp} can feed
\texttt{cmov}s that do not form a full swap (the swap move), the instruction count is not $4 \times$ comparators, and the
0-1 principle fails (SI App.~F of~\cite{alphadev}). Our method follows the same pattern (SAT encoding, symmetry breaking,
verified certificate checking) for the instruction model.
